% RASP final report, ACM SIG Proceedings format (acmart, sigconf) as required by the course.
% Build: latexmk -pdf main.tex (or upload docs/report/ to Overleaf).
% Numbers, tables and figures come from scripts/make_report_assets.py; never type a result by hand.
% Section order and word limits follow the course's report guidelines:
%   Abstract (6-8 sentences) | 1 Introduction (500-1000 words) | 2 Objective of the Work (100-200)
%   3 Major Contributions (500-1000; 3.1 papers, 3.2 experiments) | 4 Performance Measuring Metrics
%   5 Results / Outcomes (50-100 per method) | 6 Limitations and Future Scope (>= 200)
%   7 Observations from the Study (100-200). Appendices hold the detailed design and evaluation.
\documentclass[sigconf,nonacm]{acmart}
\usepackage{booktabs}
\usepackage{enumitem}
\usepackage[capitalise,noabbrev]{cleveref}
\setlist{nosep,leftmargin=*}
\setlength{\emergencystretch}{2em} % two columns and long \texttt names

% Report macros. Every result in the report is a \stat{key}: generated/*.tex (scripts/make_report_assets.py) defines
% the values with \setstat{key}{value}. A key without a value prints a red ??key and a LaTeX warning, so a number that
% has not been measured yet (Phases 4-5 before the GPU runs) is impossible to miss.
\makeatletter
\newcommand{\setstat}[2]{\expandafter\def\csname rasp@stat@#1\endcsname{#2}}
\DeclareRobustCommand{\stat}[1]{%
  \ifcsname rasp@stat@#1\endcsname\csname rasp@stat@#1\endcsname
  \else\textcolor{red}{\textbf{??#1}}\PackageWarning{rasp}{No value for `#1' (run scripts/make_report_assets.py)}\fi}
\newcommand{\stattable}[1]{\stat{table:#1}}
% \rfig[width]{name}: figures/name.pdf (or .png), or a visible placeholder until the figure exists.
\newcommand{\rfig}[2][\linewidth]{%
  \IfFileExists{figures/#2.pdf}{\includegraphics[width=#1]{figures/#2.pdf}}{%
    \IfFileExists{figures/#2.png}{\includegraphics[width=#1]{figures/#2.png}}{%
      \fbox{\parbox[c][3.5cm][c]{\dimexpr#1-2\fboxsep-2\fboxrule\relax}{\centering
        \textcolor{red}{\textbf{missing figure: \texttt{\detokenize{#2}}}}\\[2pt]
        \small (produced by \texttt{scripts/make\_report\_assets.py} after the GPU runs)}}}}}
\makeatother
\newcommand{\x}{$\times$}
\newcommand{\G}[1]{\textbf{G#1}}
\newcommand{\todo}[1]{\textcolor{orange}{[TODO: #1]}}

\InputIfFileExists{generated/platform.tex}{}{}
\InputIfFileExists{generated/phase2.tex}{}{}
\InputIfFileExists{generated/phase3.tex}{}{}
\InputIfFileExists{generated/phase4.tex}{}{}
\InputIfFileExists{generated/phase5.tex}{}{}
\InputIfFileExists{generated/omp.tex}{}{}

\settopmatter{printacmref=false}
\setcopyright{none}

\begin{document}

\title{Runtime-Adaptive Sparse Matrix Optimization for GPUs}
\subtitle{Sparse matrix representation and GPU memory-layout optimization\\[2pt]
  Course project, High Performance Computing, December 2026}

\author{Harsh Raj (S20240010084)}
\author{Karmanya (S20240010108)}
\author{Vivian (S20240010258)}
\affiliation{\institution{Indian Institute of Information Technology Sri City}\city{Sri City}\state{Andhra Pradesh}\country{India}}
\renewcommand{\shortauthors}{Harsh Raj, Karmanya and Vivian}

\begin{abstract}
Sparse matrix--vector (SpMV) and sparse matrix--matrix (SpGEMM) multiplication are memory-bound and irregular, so the
fastest GPU implementation depends on each matrix's structure and on whether the operands fit in GPU memory, yet the
two methods we build on, adaptive row-grouped CSR (ArgCSR)~\cite{heller2012argcsr} and chunked out-of-core
SpGEMM~\cite{deveci2018multilevel}, fix these choices once for all inputs.
We aim to make the choices at run time: to come within 10\% of the best of \stat{p2-configs} SpMV configurations on
every matrix without ever costing more than the vendor library for a single call, and to be at least 1.3\x{} faster
than the out-of-core baseline when a product does not fit in GPU memory.
Our library, RASP, is written in CUDA C++ on top of cuSPARSE, and we also parallelised the sequential CPU version
with OpenMP so that all CPU cores can be compared with the GPU.
Experiments run on an RTX~4050 Laptop GPU on PCIe~4.0 with an 8-core Ryzen~7 7840HS under WSL2 Ubuntu (CUDA~13.3,
cuSPARSE~12.8), with cross-checks on a Kaggle Tesla~T4 and a Core i5-1235U.
On \stat{p3-heldout-matrices} unseen SuiteSparse matrices the SpMV selector runs within \stat{p3-g2-tree}\x{} of the
per-matrix best, and on the same machine the GPU is 5.0\x{} (SpMV), 9.1\x{} (conjugate gradients) and 5.0\x{} (AMG
set-up) faster than one CPU core, and still 1.5--4.7\x{} faster than OpenMP on all eight cores, whose memory bandwidth
saturates.
Out of core, overlapping PCIe copies with computation makes SpGEMM \stat{p4-g6-AxP}\x{} ($A\cdot P$) and
\stat{p4-g6-RxA}\x{} ($R\cdot A$) faster than the baseline, short of our 1.3\x{} goal, and the data transfers still hold
it to \stat{p4-frac-of-device} of the in-memory kernel's throughput.
The code, the raw results and the scripts that regenerate every number in this report are openly available at
\url{https://github.com/Vivianmatheww/HPC_Project}.
\end{abstract}

\keywords{SpMV, SpGEMM, sparse formats, autotuning, out-of-core, CUDA, OpenMP}

\maketitle

\section{Introduction}\label{sec:intro}

Sparse matrix--vector multiplication (SpMV, $y = Ax$) and sparse matrix--matrix multiplication (SpGEMM,
$C = AB$) dominate iterative solvers, multigrid set-up and graph analytics. Both are limited by memory traffic
rather than arithmetic, and both are irregular: how evenly the nonzeros are spread over rows, how long the longest
row is, and how large the operands are relative to GPU memory all change which implementation is fastest.

The two papers this project starts from show this clearly, and both stop short of acting on it.
Heller and Oberhuber's adaptive row-grouped CSR (ArgCSR) gives long rows more threads, and its performance depends
strongly on one parameter, \texttt{desiredChunkSize}: on \texttt{rajat23} a value of~1 is six times faster than~32,
while the AFE matrices prefer~32. The paper then fixes the value to~1 for every matrix, and concludes that ``if high
performance is the top priority, one should test more formats and choose the best one''~\cite{heller2012argcsr}.
Deveci et al.\ extend the KokkosKernels SpGEMM to multilevel memory: when the operands do not fit in
fast memory they keep one operand's row block resident and stream the other through it~\cite{deveci2018multilevel}.
Their partitioning heuristic (Algorithm~4) reserves a fixed 75\%/25\% of fast memory for the resident and the
streamed operand and chooses between the two chunking variants by copy volume alone. Copies and computation are not
overlapped, which the authors leave as future work.

RASP (\emph{Runtime-Adaptive Sparse Processing}) makes these choices at run time, per matrix, from cheap
measurements:
\begin{enumerate}
\item \textbf{A lightweight SpMV selector.} One GPU kernel extracts structural features; a shallow decision tree,
  trained offline and compiled into the library as plain C++, predicts the best format \emph{and} its parameters;
  a break-even test converts formats only when the expected number of calls repays the conversion (\cref{sec:design-spmv}).
\item \textbf{An overlapped, self-calibrating out-of-core SpGEMM.} Paper~1's chunking, extended with double and
  triple buffering, a cost model that simulates the copy/compute pipeline using bandwidths and kernel rates
  measured at run time, and re-partitioning while it runs (\cref{sec:design-spgemm}).
\item \textbf{A re-evaluation of both papers} on a PCIe-attached consumer GPU (RTX~4050 Laptop, Ada) under
  WSL2, against the 2026 cuSPARSE (\cref{sec:eval}).
\end{enumerate}
Automatic format selection with machine learning is an established research area
(\cref{sec:related}); our novelty is relative to the two assigned papers: tuning ArgCSR's parameters per matrix and
per row region, a conversion-aware cost model, and bringing runtime adaptivity and overlap to out-of-core SpGEMM.

\Cref{tab:goals} lists the goals we set at the start and where each is evaluated.

\begin{table}[t]
\centering\small
\caption{Project goals and where they are evaluated.}\label{tab:goals}
\begin{tabular}{@{}lp{6.2cm}p{5.6cm}l@{}}
\toprule
 & Goal & Target & Section \\
\midrule
\G1 & Reproduce paper 2 (ArgCSR) on modern hardware & same qualitative winners and losers & \ref{sec:eval-g1} \\
\G2 & Runtime SpMV selector & geomean $T_\text{selected}/T_\text{oracle} \le 1.10$; never $>$5\% slower than cuSPARSE on matrices $\ge$ 1M nnz & \ref{sec:eval-g2} \\
\G3 & Selection is cheap & features + inference $\le$ 2 SpMVs & \ref{sec:eval-g2} \\
\G4 & Conversion-aware choice & never $>$5\% worse than cuSPARSE at $N=1$ & \ref{sec:eval-g2} \\
\G5 & Paper 1's placement study on a PCIe GPU & our version of its Table 3 & \ref{sec:eval-g5} \\
\G6 & Overlapped, adaptive out-of-core SpGEMM & $\ge 1.3\times$ over Algorithm 4 when the problem exceeds the budget & \ref{sec:eval-g6} \\
\G7 & End-to-end applications & CG (SpMV) and AMG set-up (SpGEMM) working and timed & \ref{sec:eval-g7} \\
\bottomrule
\end{tabular}
\end{table}

\section{Objective of the Work}\label{sec:objective}

The objective is to replace the static choices of the two papers with choices made at run time, and to show on
measurements that this is fast, cheap and safe. For SpMV, a selector should pick a format and its parameters for
each matrix, come within 10\% of the best of all configurations on matrices it has never seen, cost no more than about
two SpMVs to decide, and never be slower than cuSPARSE when the matrix is used only once. For SpGEMM, a planner should
keep every product within a fixed GPU memory budget, overlap PCIe copies with computation, and be at least 1.3\x{}
faster than paper~1's Algorithm~4 whenever the problem exceeds the budget. Along the way we re-test both papers'
claims on current hardware and libraries, and we show both techniques at work inside real applications (CG and AMG
set-up). Finally, we compare the GPU with sequential and OpenMP CPU versions on the same machine.
\Cref{tab:goals} lists the measurable goals and where each is evaluated.

\begin{table*}[t]
\centering\small
\caption{Project goals and where they are evaluated.}\label{tab:goals}
\begin{tabular}{@{}lp{6.4cm}p{6.6cm}l@{}}
\toprule
 & Goal & Target & Section \\
\midrule
\G1 & Reproduce paper 2 (ArgCSR) on modern hardware & same qualitative winners and losers & \ref{sec:res-g1} \\
\G2 & Runtime SpMV selector & geomean $T_\text{selected}/T_\text{oracle} \le 1.10$; never $>$5\% slower than cuSPARSE on matrices $\ge$ 1M nnz & \ref{sec:res-g2} \\
\G3 & Selection is cheap & features + inference $\le$ 2 SpMVs & \ref{sec:res-g2} \\
\G4 & Conversion-aware choice & never $>$5\% worse than cuSPARSE at $N=1$ call & \ref{sec:res-g4} \\
\G5 & Paper 1's placement study on a PCIe GPU & our version of its Table 3 & \ref{sec:res-g5} \\
\G6 & Overlapped, adaptive out-of-core SpGEMM & $\ge 1.3\times$ over Algorithm 4 when the problem exceeds the budget & \ref{sec:res-g6} \\
\G7 & End-to-end applications & CG (SpMV) and AMG set-up (SpGEMM) working and timed & \ref{sec:res-cg}, \ref{sec:res-amg} \\
\bottomrule
\end{tabular}
\end{table*}

\section{Major Contributions}\label{sec:contrib}

This project contributes four things, described in full in \cref{sec:design}:
\begin{enumerate}
\item \textbf{A lightweight runtime SpMV selector.} One GPU kernel extracts structural features; a shallow decision
  tree, trained offline and compiled into the library as plain C++, chooses among \stat{p2-configs} configurations of
  CSR, ELL, SELL, ArgCSR, cuSPARSE and two new kernels (ArgCSR with a chunk size per row region, and CSR with lanes
  per row); a break-even test converts formats only when the expected number of calls repays the conversion.
\item \textbf{An overlapped, self-calibrating out-of-core SpGEMM.} Paper~1's chunking, extended with double and
  triple buffering, a cost model that simulates the copy/compute pipeline from bandwidths and kernel rates measured at
  run time, and re-partitioning while it runs, all within a hard device-memory budget.
\item \textbf{A re-evaluation of both papers} on a PCIe-attached consumer GPU against the 2026 cuSPARSE.
\item \textbf{Sequential and OpenMP CPU versions} of the same kernels, for same-machine comparisons.
\end{enumerate}

\subsection{Contributions of the papers we studied}\label{sec:papers}
\paragraph{Paper 1: Deveci et al., multilevel-memory SpGEMM~\cite{deveci2018multilevel}.} The paper runs the
KokkosKernels \texttt{kkmem} SpGEMM~\cite{deveci2018kkmem} on a KNL and on a P100 connected to its host over NVLink.
It contributes (i) two chunked out-of-core algorithms: \emph{chunk1} keeps a row block of $A$ and $C$ resident and
streams the row blocks of $B$, \emph{chunk2} keeps a row block of $B$ resident and streams $A$ and $C$; (ii) a
partitioning heuristic, Algorithm~4, that splits fast memory 75\%/25\% between the resident and the streamed operand
and picks the variant with the smaller copy volume; and (iii) a placement study (its Table~3) showing that $B$ in host
memory costs 7--29\x, while $A$ and $C$ hurt only when large, and that unified memory reaches at most 30\% of the
device's speed.

\paragraph{Paper 2: Heller and Oberhuber, adaptive row-grouped CSR~\cite{heller2012argcsr}.} The paper contributes
the ArgCSR format: rows are packed into groups handled by one CUDA block, spare threads go to the rows with the
longest chunks, and chunks are padded and interleaved so that reads coalesce. Measured on 1,600 matrices on a Fermi
C2070, ArgCSR beat CPU CSR on 1,168 of them and was up to ten times faster than the cuSPARSE of 2012 on circuit
matrices. The paper also documents how strongly its chunk-size parameter matters, and recommends testing several
formats per matrix without automating that choice.

\paragraph{Other work we drew on.} CSR-vector and ELL kernels~\cite{bell2009spmv}, SELL-$C$-$\sigma$~\cite{kreutzer2014sell},
CSR5~\cite{liu2015csr5} and merge-based SpMV~\cite{merrill2016merge} define the candidate formats. SMAT~\cite{li2013smat},
decision trees~\cite{sedaghati2015selection}, SVMs~\cite{benatia2016svm} and CNNs~\cite{zhao2018cnn} have all been used to
select formats; unlike them, our selector also predicts the cost of \emph{acting} on its choice. Hash-based SpGEMM
(nsparse~\cite{nagasaka2017nsparse}, spECK~\cite{parger2020speck}) and output-size handling~\cite{liu2014spgemm} inform the
SpGEMM engine.

\subsection{Experiments: test-bed and datasets}\label{sec:testbed}
\paragraph{Hardware and software.} The primary machine is a laptop with an RTX~4050 Laptop GPU (Ada, sm\_89, 6\,GB,
20~SMs, 24\,MB L2) on PCIe~4.0~x8 and an 8-core, 16-thread Ryzen~7 7840HS, running Ubuntu~22.04 under WSL2 with CUDA~13.3,
cuSPARSE~12.8 and g++~11.4. \Cref{tab:platform} gives its measured characteristics. A Kaggle Tesla~T4 (CUDA~12.8, 4-thread
Xeon) repeats the same-machine comparison, and the OpenMP scaling study also ran on a Core i5-1235U. No simulator is
used: everything is measured on real hardware.

\paragraph{Datasets.} SpMV: \stat{p2-matrices} SuiteSparse matrices~\cite{davis2011suitesparse} for training (50 drawn at
random, real and square with $10^4$--$5\cdot10^6$ nonzeros, plus 10 named in paper~2) and \stat{p3-heldout-matrices}
different matrices drawn the same way for evaluation, never seen in training. CG: the SPD matrices of these sets, 20 more
SPD matrices, and 3-D Laplacians. SpGEMM: paper~1's multigrid products on the 7-point Laplace3D with $2\times2\times2$
aggregation ($P$, $R=P^T$), generated in memory, which allows multi-GB problems.

\paragraph{Procedure.} Kernels are timed with CUDA events: 5 warm-up runs, then the median of 20 with the
interquartile range; the L2 cache is flushed before every timed SpMV. Out-of-core SpGEMM is timed from pinned inputs to
pinned output, including every transfer and the planning time. Methods are run in a freshly shuffled order in every
repetition, because a fixed order produced fake differences of $\pm9\%$ under WSL2. Every result is checked against a
CPU reference, and every CSV records the GPU, driver, library versions and git commit. Runs use AC power and the
``best performance'' mode. \texttt{scripts/run\_all.sh} rebuilds from a clean tree, reruns every test and benchmark,
and regenerates every number, table and figure in this report.

\begin{table}[t]
\centering\small
\caption{Platform measurements of the RTX~4050 laptop (\texttt{results/platform/uvm\_probe.csv}).}\label{tab:platform}
\resizebox{\columnwidth}{!}{%
\begin{tabular}{@{}lr@{}}
\toprule
Device memory, sequential read (4 GiB) & \stat{pf-device-read} GB/s \\
Device memory, copy (read + write) & \stat{pf-device-copy} GB/s \\
Pinned host $\to$ device / device $\to$ host & \stat{pf-h2d} / \stat{pf-d2h} GB/s \\
Both directions at once (one copy engine) & \stat{pf-bidir} GB/s \\
Zero-copy / managed (UVM) read by a kernel & \stat{pf-zc-read} / \stat{pf-uvm-read} GB/s \\
Random 8-byte gathers, device vs zero-copy & \stat{pf-gather-device} vs \stat{pf-gather-zc} M/s \\
\bottomrule
\end{tabular}}
\end{table}

\section{Performance Measuring Metrics}\label{sec:metrics}

\paragraph{Throughput (GFLOP/s).} SpMV performs $2\,\mathrm{nnz}$ flops per call and SpGEMM two flops per
multiply-add. GFLOP/s normalises time by the work done, so matrices of different sizes can be compared, and it is the
unit both papers report, which lets us compare with them directly.

\paragraph{Effective bandwidth (fraction of measured peak).} SpMV and SpGEMM are memory-bound, so the fraction of the
\emph{measured} \stat{pf-device-read}\,GB/s a kernel reaches says how close it is to the hardware limit, and therefore
how much any other format could still gain. We use the measured rather than the datasheet peak (192\,GB/s), which no
kernel can reach. Nsight Compute's DRAM, L2 and occupancy counters explain \emph{why} a kernel sits where it does.

\paragraph{Slowdown against the oracle, $T/T_\text{oracle}$.} The oracle is the fastest of all configurations for
each matrix, found by exhaustive measurement. The ratio says how much of the attainable speed a selector leaves
behind, independently of how good cuSPARSE happens to be on that matrix. We aggregate ratios with the \emph{geometric}
mean, the correct mean for ratios (a 2\x{} gain and a 2\x{} loss cancel), and also report the worst case, since a
selector that is good on average but occasionally very slow is not safe to deploy.

\paragraph{Performance profiles.} Dolan--Mor\'e profiles~\cite{dolan2002profiles} plot, for every factor $\tau$, the
fraction of matrices on which a method is within $\tau$ of the oracle. They show the whole distribution, not only its
mean, and are the standard way to compare solvers over a test set.

\paragraph{Selection overhead in SpMV equivalents.} Feature extraction and conversion times are divided by the time of
one cuSPARSE SpMV on the same matrix. This expresses an overhead in the unit that matters to the caller: how many calls
it takes to pay it back.

\paragraph{Total time for $N$ calls.} Plan cost plus $N$ SpMVs, relative to always using cuSPARSE, for $N$ from 1 to
1000. This is the quantity break-even optimises, and it makes the one-call case ($N=1$) explicit.

\paragraph{Speedup.} Against Algorithm~4 for out-of-core SpGEMM (the paper's own baseline), against cuSPARSE for the
selector and CG, and against the sequential CPU code on the same machine for the GPU and OpenMP versions. For
OpenMP we also report speedup over the same code on one thread and parallel efficiency, which separate parallel
scaling from compiler or algorithm effects.

\paragraph{Data transfer cost.} PCIe bandwidth (host to device, device to host, both), the slowdown when an operand
lives in host memory, and out-of-core throughput as a fraction of the device-only kernel's. These quantify how much the
host link costs, which is the central constraint of out-of-core work.

\paragraph{Time to solution.} For CG and the AMG set-up we time the whole application, including planning and
transfers, because a faster kernel only matters if it shortens the end result.

\paragraph{Cost-model accuracy.} The ratio of the planner's predicted execution time to the measured one shows whether
the simulated pipeline is trustworthy enough to choose plans with.

\section{Results / Outcomes}\label{sec:results}

Each method is summarised below; \cref{sec:eval} gives the full tables and figures.

\subsection{Reproducing ArgCSR (G1)}\label{sec:res-g1}
The chunk-size sensitivity reproduces. On \texttt{rajat23}, ArgCSR with chunk size~1 runs at
\stat{p2-rajat23-d1}\,GFLOP/s and 32 at \stat{p2-rajat23-d32} (paper~2: 5.1 vs 0.81), because one long row pads its
whole group. No global value is right: fixing the paper's choice of~1 costs \stat{p2-d1-slowdown}\x{} on average against
the best value per matrix. The paper's headline does not hold against 2026 cuSPARSE: ArgCSR reaches
\stat{p2-argcsr-default-fp64}\x{} of cuSPARSE's speed in fp64 at its default and \stat{p2-argcsr-tuned-fp64}\x{} when
tuned per matrix (\cref{sec:eval-g1}). G1 is met.

\subsection{SpMV formats and our two kernels}\label{sec:res-formats}
The fastest configuration per matrix is \stat{p2-oracle-fp64}\x{} faster than cuSPARSE \texttt{ALG1} in fp64
(geometric mean), and the headroom shrinks with size: \stat{p2-oracle-size-lt100k}\x{} below $10^5$ nonzeros,
\stat{p2-oracle-size-gt1M}\x{} above $10^6$. cuSPARSE is fastest on \stat{p2-cusparse-wins-fp64} matrices and CSR-vector
with few lanes on most others. Without tuning, our ArgCSR-local reaches \stat{p2-local-vs-tuned} of per-matrix-tuned
ArgCSR, and CSR-binned matches tuned CSR-vector, but neither moves the oracle by more than about a percent.

\subsection{Runtime SpMV selector (G2, G3)}\label{sec:res-g2}
On \stat{p3-heldout-matrices} unseen matrices the decision tree runs at \stat{p3-g2-tree}\x{} of the oracle (target
1.10: met), \stat{p3-g2-tree-vs-cusparse}\x{} faster than always using cuSPARSE, which sits at \stat{p3-g2-cusparse}\x;
hand-written rules reach only \stat{p3-g2-rules}\x. It is never more than 5\% slower than cuSPARSE on matrices with at
least $10^6$ nonzeros, though only \stat{p3-heldout-large} held-out samples are that large. Deciding costs a median of
\stat{p3-g3-median} SpMVs, so G3 is met only narrowly: \stat{p3-g3-size-gt1M} on large matrices but
\stat{p3-g3-size-lt100k} on small ones.

\subsection{Conversion-aware break-even (G4)}\label{sec:res-g4}
With break-even, planning plus a single call costs \stat{p3-g4-tree-1}\x{} the time of always using cuSPARSE, and
\stat{p3-g4-n1-same} of \stat{p3-g4-n1-total} one-call plans do exactly the baseline's work. Without break-even the same
tree costs \stat{p3-g4-tree-nobe-1}\x{} at one call (worst \stat{p3-g4-tree-nobe-1-worst}\x), because it converts formats
it cannot repay. At 1000 calls the selector gains \stat{p3-g4-tree-1000}\x. G4 is met: the selector never costs more than
it saves when a matrix is used once.

\subsection{Operand placement (G5)}\label{sec:res-g5}
With $B$, the operand every row gathers from at random, in pinned host memory, cuSPARSE SpGEMM slows down
\stat{p4-place-RxA-b}\x{} ($R\cdot A$) and \stat{p4-place-AxP-b}\x{} ($A\cdot P$); the streamed left operand of
$R\cdot A$ costs only \stat{p4-place-RxA-a}\x. Managed memory behaves exactly like pinned memory under WSL2
(\stat{p4-managed-vs-pinned}\x). Our ratios are smaller than paper~1's 7--29\x{} on NVLink because our device-resident
kernel is slower, so there is less speed to lose, but the conclusion is the same: $B$ stays on the device and $A$
and $C$ are streamed.

\subsection{Overlapped out-of-core SpGEMM (G6)}\label{sec:res-g6}
When the problem exceeds the budget, RASP is \stat{p4-g6-AxP}\x{} faster than Algorithm~4 on $A\cdot P$ (G6
\stat{p4-g6-AxP-verdict}) but only \stat{p4-g6-RxA}\x{} on $R\cdot A$ (\stat{p4-g6-RxA-verdict}). Most of the gain is
copy/compute overlap, not the cost model: double buffering with Algorithm~4's own partitions already gives
\stat{p4-speedup-alg4-overlap-AxP}\x{} and \stat{p4-speedup-alg4-overlap-RxA}\x. On $R\cdot A$, RASP cuts $A$ and $C$
into too many tiles once $B$ is split, and re-planning never fires in those plans. Even overlapped, PCIe holds
throughput to \stat{p4-frac-of-device} of the device-only kernel.

\subsection{Conjugate gradients (G7)}\label{sec:res-cg}
Jacobi-preconditioned CG on \stat{p5-cg-matrices} SPD matrices, using the solver's iteration cap as the call count, is
\stat{p5-cg-tree-geomean}\x{} faster than the same solver with cuSPARSE (worst \stat{p5-cg-tree-worst}\x, best
\stat{p5-cg-tree-best}\x), and \stat{p5-cg-failures} solves fail the CPU residual check. The gain is small because SpMV
is only \stat{p5-cg-spmv-share}\% of a solve. An overstated hint does no harm: the cap and the exact iteration count
pick the same kernel on \stat{p5-cg-hint-same} matrices, and planning costs a median \stat{p5-cg-plan-pct-median}\% of a
solve.

\subsection{AMG set-up (G7)}\label{sec:res-amg}
Building a Laplace3D ($n=192$) hierarchy under a 256\,MiB budget, RASP is \stat{p5-amg-192-rasp-speedup}\x{} faster than
Algorithm~4 and about as fast as overlapped Algorithm~4 (\stat{p5-amg-192-alg4-overlap-speedup}\x). Only the finest
levels are tiled, and level~0 takes 82--85\% of the set-up time. Every method produces identical coarse operators that
satisfy the Galerkin row-sum identity. The \stat{p5-amg-256-rasp-speedup}\x{} measured at $n=256$ is inflated by an
anomalous host-side estimate in that Algorithm~4 run (\cref{sec:eval-g7}).

\subsection{GPU and OpenMP against one CPU core}\label{sec:res-seq}
\Cref{tab:same-machine} compares all three versions on the same machines. On the RTX~4050 laptop the GPU is
5.0\x{} faster than one core for SpMV, 9.1\x{} for CG and 5.0\x{} for the AMG set-up, including PCIe transfers. OpenMP
on the eight Ryzen cores reaches 3.3\x, 1.9\x{} and 3.0\x, so the GPU's lead over the whole CPU is 1.5\x, 4.7\x{} and
1.7\x. The T4's ratios are higher mainly because Kaggle's 2\,GHz Xeon is slower. With a warm L2 cache, the T4's
best-configuration lead over cuSPARSE falls from 1.44\x{} to 1.28\x, so we report flushed numbers.

\begin{table}[t]
\centering\small
\caption{Speedup over one core of the same CPU (\texttt{results/*/comparison.txt}). OpenMP: best thread count.
SpMV: cuSPARSE, geometric mean over six matrices, fp64, L2 flushed. CG: Laplace $128^3$. AMG: set-up at $n=160$, in
core, GPU including transfers.}\label{tab:same-machine}
\resizebox{\columnwidth}{!}{%
\begin{tabular}{@{}lrrrr@{}}
\toprule
 & \multicolumn{2}{c}{RTX 4050 / Ryzen 7 7840HS} & \multicolumn{2}{c}{Tesla T4 / Xeon 2\,GHz} \\
\cmidrule(lr){2-3}\cmidrule(l){4-5}
 & OpenMP & GPU & OpenMP & GPU \\
\midrule
SpMV & 3.3\x{} (8 t) & 5.0\x & 2.3\x{} (4 t) & 10.4\x \\
CG, time to solution & 1.9\x{} (4 t) & 9.1\x & 2.7\x{} (4 t) & 24.0\x \\
AMG set-up & 3.0\x{} (16 t) & 5.0\x & 1.6\x{} (2 t) & 4.8\x \\
\bottomrule
\end{tabular}}
\end{table}

\subsection{OpenMP multi-core version}\label{sec:res-omp}
Parallelising the sequential code with OpenMP, without changing its algorithms and verified against it, gives at most
\stat{omp-spmv-best-speedup}\x{} for SpMV, \stat{omp-cg-best-speedup}\x{} for CG and \stat{omp-amg-best-speedup}\x{} for
the AMG SpGEMMs on the 12-thread i5-1235U. SpMV and CG are capped by memory bandwidth, which rises only
\stat{omp-bw-ceiling}\x{} from one thread to all; the SpGEMM is capped by single-threaded zero-filling and the page faults
it causes (Amdahl's law). The GPU's advantage on these kernels is bandwidth, which no number of CPU cores supplies
(\cref{sec:omp}).

\subsection{Kernel profiles}\label{sec:res-ncu}
Nsight Compute confirms the bandwidth explanation. CSR-vector on \texttt{af\_shell1} keeps DRAM busy
\stat{p5-ncu-spmv-csr-vector-af-shell1-dram}\% of the time at \stat{p5-ncu-spmv-csr-vector-af-shell1-gbs}\,GB/s, and
cuSPARSE sits at the same limit, so no format can gain much there. ArgCSR with chunk size 32 on \texttt{rajat23} moves
only 1.3\x{} the DRAM bytes of chunk size 1 but takes 17\x{} as long: its warps stall on barriers at
\stat{p5-ncu-spmv-argcsr-D32-rajat23-occ}\% occupancy, so its cost is synchronisation, not padding traffic.

\section{Limitations}\label{sec:limitations}
\begin{itemize}
\item \textbf{One laptop GPU under WSL2.} Every number comes from one RTX~4050 Laptop. Clocks and thermals vary, which
  we contain with medians, shuffled orders and AC power but cannot remove. WSL2 also changes the semantics we
  study: managed memory never migrates, so paper~1's UVM experiments cannot be reproduced here, and host-side
  latencies of tens of microseconds weigh on small SpMVs and on G3. A native Linux system and a data-centre GPU would
  test whether the trained selector transfers across architectures.
\item \textbf{Few large held-out matrices.} Only \stat{p3-heldout-large} held-out samples have at least $10^6$
  nonzeros, so the ``never lose on large matrices'' half of G2 rests mostly on cross-validation.
\item \textbf{G3 is borderline} (\stat{p3-g3-median} SpMVs median): fine for large matrices, expensive for small
  ones, where the gate simply skips selection.
\item \textbf{SpGEMM engine.} cuSPARSE's SpGEMM does not report its workspace in advance and blocks the host in its
  size queries. We learn the workspace size, plan pipeline depths around the host stalls, and split tiles that do
  not fit. A custom hash SpGEMM with a column-range filter (as \texttt{kkmem}) would avoid the full-height row pointers
  of $B$ blocks and the host round trips. Our partial-product sum uses one thread per row, which suits multigrid
  products but not power-law graphs.
\item \textbf{Workloads.} Out-of-core results are for paper~1's multigrid products on Laplace3D. Graph products
  ($L\cdot L$ for triangle counting) are supported by the benchmark but not part of the evaluated set.
  \todo{remove this item if the SNAP runs are added.}
\item \textbf{Budget emulation.} The fast-memory budget emulates a smaller GPU, as paper~1 does; the operands live in
  pinned host memory and the estimation pass runs on the host CPU.
\end{itemize}

\section{Conclusion}\label{sec:conclusion}
Both source papers measured the right thing and then fixed a choice. We made the choices at run time. For SpMV, a
GPU feature pass, a small cross-validated tree and a break-even test reach \stat{p3-g2-tree}\x{} of the per-matrix
oracle on unseen matrices while never costing more than cuSPARSE for a single call; on the way, ArgCSR's sensitivity
to its chunk size reproduced while its advantage over the 2026 vendor library did not. For out-of-core SpGEMM, a
planner that measures the PCIe link, simulates the copy/compute pipeline and re-plans from a probe block is
\stat{p4-g6-AxP}\x{} and \stat{p4-g6-RxA}\x{} faster than paper~1's Algorithm~4 on the multigrid products.
\todo{final sentence once Phase 4--5 results are in.}

\paragraph{Reproducibility.} \texttt{scripts/run\_all.sh} rebuilds from a clean tree, runs every test and benchmark
and regenerates every number, table and figure in this report (\texttt{scripts/make\_report\_assets.py}). Raw CSVs,
with environment stamps, are in \texttt{results/}.

\section{Observations from the Study}\label{sec:observations}

The study asked whether the choices the two papers fixed statically can be made at run time, cheaply and safely.
It was designed so that each answer is measured against a fair baseline: an exhaustive oracle, held-out matrices,
the papers' own algorithms, and CPU versions on the same machine. For SpMV the answer is yes: a GPU feature pass, a
small tree and a break-even test reach \stat{p3-g2-tree}\x{} of the oracle on unseen matrices and never cost more
than cuSPARSE for a single call. For out-of-core SpGEMM the answer is only partly yes: overlapping copies with
computation is what pays, and a cost model is worth having only if it does not over-partition. End to end, a faster
kernel helps only as much as its share of the application, which is why CG gains \stat{p5-cg-tree-geomean}\x{} while
the GPU itself is 9\x{} faster than a CPU core. Two questions remain open: whether the selector transfers to other
GPUs without retraining, and how out-of-core SpGEMM behaves with unified-memory oversubscription on native Linux.


\bibliographystyle{ACM-Reference-Format}
\bibliography{refs}

\appendix
\section{Design}\label{sec:design}

RASP is a CUDA C++ library (\texttt{include/rasp/}, \texttt{src/}) with two runtime entry points: \texttt{SpmvPlan}
for SpMV and \texttt{SpgemmContext} for out-of-core SpGEMM. Both decide once per matrix (SpGEMM also while running),
explain their decision in a report object, and fall back to a safe choice whenever a prediction turns out wrong.

\subsection{SpMV: candidates}\label{sec:design-spmv}
The candidate set has \stat{p2-configs} configurations behind one interface: cuSPARSE \texttt{CSR\_ALG1} and
\texttt{ALG2}; CSR-scalar; CSR-vector with $v \in \{2,\dots,32\}$; ELL; SELL-$C$-$\sigma$ with $C\in\{8,32\}$ and
$\sigma \in \{1, 256, n\}$; ArgCSR with $\mathtt{blockDim} \in \{64,128,256\}$ and
$\mathtt{desiredChunkSize} \in \{1,\dots,32\}$, reimplemented from the paper's Listing~2 and checked against its
Figure~3; and two extensions aimed at the gaps the Phase~2 sweep exposed.
\begin{itemize}
\item \textbf{ArgCSR-local} chooses \texttt{desiredChunkSize} per window of 4,096 rows instead of globally. The sweep
  showed that the best global value tracks the mean row length and keeps a few dozen rows per group, so each window
  gets $D = \mathrm{pow2round}(R\cdot\bar{\ell}/B)$, halved when the window's row-length coefficient of variation
  exceeds~1 and clamped to $[1, 32]$. The kernel is unchanged, since groups already carry their own chunk size.
\item \textbf{CSR-binned} is CSR-vector with the number of lanes chosen per row,
  $v = \mathrm{pow2ceil}(\lceil \ell/t \rceil)$ up to 32, and a whole block for rows longer than 4,096. Rows are grouped
  by bin and one launch covers all bins.
\end{itemize}
Every format has a GPU converter that produces arrays bit-identical to the host converter (tested). Moving
conversion to the GPU lowered its median cost from tens of SpMV calls to \stat{p3-gpu-convert-calls-ell} (ELL),
\stat{p3-gpu-convert-calls-sell} (SELL), \stat{p3-gpu-convert-calls-argcsr} (ArgCSR) and
\stat{p3-gpu-convert-calls-csr-binned} (CSR-binned) cuSPARSE SpMVs, which is what makes a break-even decision meaningful.

\subsection{SpMV: features, selector and break-even}
A single feature kernel reads the row pointers plus the first and last column index of each row (far less traffic
than one SpMV) and accumulates row-length mean, coefficient of variation and maximum, the fraction of empty rows,
ELL and SELL-32 padding factors, the share of nonzeros in rows longer than 4,096, the mean column span, and an
eight-bucket row-length histogram.

The selector went through two versions. Hand-written thresholds derived from the Phase~2 data (v1) were no better
than always using cuSPARSE on unseen matrices (\cref{sec:eval-g2}). The shipped selector (v2) is trained by
\texttt{scripts/train\_selector.py} on \stat{p3-train-samples} samples (\stat{p3-train-matrices} matrices in fp32 and
fp64 from \stat{p3-train-groups} SuiteSparse groups). It first picks a small \emph{portfolio} greedily: five
configurations (cuSPARSE \texttt{ALG1}, CSR-vector $v{=}4$, CSR-binned $t{=}2$, CSR-scalar, cuSPARSE \texttt{ALG2})
whose per-matrix best is within \stat{p3-cv-portfolio}\x{} of the full \stat{p2-configs}-configuration oracle. A
decision tree then chooses among them, validated by 5-fold cross-validation \emph{grouped by SuiteSparse group}, so
matrix families never straddle training and test. Cross-validation chose depth~\stat{p3-tree-depth}
(\stat{p3-tree-leaves} leaves, splitting on the logarithms of the maximum row length and of nnz) and a large-matrix
gate that always uses cuSPARSE \texttt{ALG1} above 1M nonzeros: the gate improved every metric. The tree is emitted as
plain C++ (\texttt{src/selector/tree\_generated.hpp}) together with the cost models fitted in the same run, so there
is no runtime ML dependency.

\texttt{SpmvPlan(A, x, y, \{expected\_calls = N\})} decides in four steps:
\begin{enumerate}
\item \emph{Feature gate.} Features are extracted only if
  $N \cdot \hat T_\text{cuSPARSE} \cdot g > \hat T_\text{features}$, where $\hat T$ are fitted cost models (a launch
  cost plus bytes over an effective bandwidth) and $g = \stat{p3-expected-gain-pct}\%$ is the tree's
  \emph{cross-validated expected} saving over cuSPARSE. An earlier version used the maximum training gain
  and extracted features far too eagerly.
\item \emph{Steady-state choice} by the tree.
\item \emph{Break-even.} If the chosen format needs a conversion, it is kept only if
  $N\cdot\hat T_\text{cuSPARSE}(1 - 1/\hat s) > \hat T_\text{convert}$; otherwise the best choice that needs no
  conversion is used.
\item \emph{Build} with the GPU converter. A format refused for excessive padding falls back to the no-conversion choice.
\end{enumerate}
A report states the configuration, the steady-state choice, the measured feature and conversion times and the reason.

\subsection{Out-of-core SpGEMM}\label{sec:design-spgemm}
\texttt{SpgemmContext ctx(budget)} computes $C = AB$ for $A$, $B$, $C$ in pinned host memory using at most
\texttt{budget} bytes of device memory, as paper~1 emulates smaller GPUs with fixed budgets. The context allocates
the whole budget once and carves every device buffer from it with a first-fit range allocator, so the budget holds
exactly and the peak is recorded. It also creates a copy stream and a compute stream and measures pinned PCIe
bandwidth with 64\,MiB copies. cuSPARSE's generic SpGEMM is the in-core engine. A small core function runs it on
non-owning views with every buffer (both workspaces, and $C$'s arrays once $\mathrm{nnz}(C)$ is known) supplied by
the caller, so that buffers can live in the arena, or in pinned or managed memory for the placement study.

\paragraph{Estimation.} One host pass over $A$'s column indices computes $f_i$ and the upper bound
$\min(f_i, n_B)$ on each row of $C$, with parallel prefix sums. Both planners use the same estimate.

\paragraph{Execution modes.} \emph{In-core} (everything fits); \emph{zero-copy} ($B$ in the budget, $A$ read and
$C$ written in place over PCIe, row chunks bounding the workspace); \emph{ChunkB} (paper~1's chunk2: a row block of
$B$ resident, $A$/$C$ row blocks streamed; with one $B$ block this is the common 1-D case); and \emph{ChunkAC}
(paper~1's chunk1: an $A$/$C$ row block resident, $B$ row blocks streamed). As in paper~1, $A$ is not split by
column. A block of $B$ is presented to cuSPARSE as a full-height matrix whose rows outside the block are empty, so
$A$'s other columns contribute nothing. Partial products are summed on the device by our own two-pass merge
(count, scan, fill). With several $B$ blocks, ChunkB parks partial $C$ blocks in pinned memory between passes.

\paragraph{Pipeline.} Every mode is a list of tiles run by one loop. Tile $k+d-1$ is uploaded on the copy stream
before tile $k$'s cuSPARSE calls, because those calls block the host in their size queries; downloads queue on the
same stream, matching the single copy engine. Buffers return to the arena through events, so no stream can touch a
buffer after it is reused. The planner chooses the depth $d$; $d=1$ is strictly sequential.

\paragraph{Planning by simulation.} For every feasible mode, \texttt{plan\_rasp} gives the streamed operand
\emph{all} the memory the resident operand leaves free, rather than a fixed 25\%, and searches block sizes and depths
$d\in\{2,3\}$. Each candidate is scored by simulating the executor's issue order on one copy engine: per tile, an
upload at the measured H2D bandwidth, a computation at $f/R$ plus a per-call overhead, and a download. The memory
model charges each mode with exactly the buffers the executor holds at once (for 1-D ChunkB, $d$ blocks each of
$A$ and $C$, the resident $B$ and cuSPARSE's workspace, estimated as $\omega f$). $C$'s size is estimated as
$\gamma$ times the upper bound, with $\gamma = 1$ until measured. The simulation also covers in-core execution and
picks overlapped chunking when that is faster even though everything fits.

\paragraph{Calibration and re-planning.} For the common 1-D case the first block is a small \emph{probe} holding
about an eighth of a block's work. Nothing overlaps the pipeline fill anyway, so the probe is nearly free. From it
and every later tile the executor measures the kernel rate $R$, the compression $\gamma = \mathrm{nnz}(C)/$bound and
the workspace per multiply-add $\omega$. When any of them drifts more than 20\% from what the plan assumed, the
blocks not yet uploaded are re-partitioned with the measured values (with a 10\% safety margin).

\paragraph{Correctness never depends on the estimates.} A tile whose workspace or output does not fit, even after
waiting for every pending free, is split in half and redone. 1-D tiles split over $A$ rows; tiles that sum partial
products split over $B$ rows. Only a single row that cannot fit raises \texttt{BudgetTooSmall}.

\paragraph{Baseline: Algorithm 4.} \texttt{plan\_alg4} reimplements paper~1's heuristic: in-core if everything fits,
otherwise 75\%/25\% of the budget for the resident/streamed operand, binary-search partitions, and the variant with
the smaller copy volume, run at depth~1. Partial-$C$ buffers are charged to the $A$/$C$ share and the cuSPARSE
workspace to the streamed share. A depth-2 version isolates the effect of overlap alone.

\section{Evaluation}\label{sec:eval}

\subsection{Baseline: how well does cuSPARSE use the GPU?}
cuSPARSE \texttt{CSR\_ALG1} reaches a median of \stat{p1-fp64-csr-alg1-bw}\% of the measured device bandwidth in
fp64 (best \stat{p1-fp64-csr-alg1-bw-max}\%) over the \stat{p1-matrices} training matrices, and is faster than
\texttt{ALG2} on \stat{p1-alg1-wins-fp64} of them. Size dominates: the median is \stat{p1-fp64-bw-lt100k}\% below
$10^5$ nonzeros, \stat{p1-fp64-bw-100k-1M}\% up to $10^6$ and \stat{p1-fp64-bw-gt1M}\% above, because a small SpMV
lasts a few microseconds and is limited by launch latency. Any headroom for format selection therefore lies mostly
in small and medium matrices.

\subsection{G1: reproducing paper 2 (ArgCSR)}\label{sec:eval-g1}

\begin{figure}[t]
\centering\rfig[0.62\linewidth]{p2_argcsr_chunk}
\caption{ArgCSR (fp64, \texttt{blockDim} 128) against \texttt{desiredChunkSize}: the two matrices move in opposite
directions, as in paper~2.}\label{fig:chunk}
\end{figure}

\begin{table}[t]
\centering\small
\caption{Paper~2's named matrices (fp64, GFLOP/s): cuSPARSE \texttt{ALG1}, ArgCSR at the paper's default and tuned,
and the fastest configuration.}\label{tab:named}
\stattable{p2-named}
\end{table}

\textbf{The chunk-size sensitivity reproduces.} On \texttt{rajat23}, ArgCSR with \texttt{desiredChunkSize} 1 runs at
\stat{p2-rajat23-d1}\,GFLOP/s and with 32 at \stat{p2-rajat23-d32} (paper~2: 5.1 vs 0.81 on a Fermi C2070). Padding
causes the collapse: one long row forces its whole group to the long chunk, and the stored slots grow from
\stat{p2-rajat23-stored-d1}\x{} to \stat{p2-rajat23-stored-d32}\x{} the number of nonzeros. \texttt{af\_shell1} goes the other way
(\stat{p2-afshell1-d32} vs \stat{p2-afshell1-d1}\,GFLOP/s for 32 vs 1), as paper~2 reported for the AFE set
(\cref{fig:chunk}).

\textbf{But no global value is right, and paper~2's ``robust'' choice is rarely the best.} \texttt{desiredChunkSize}~1
is the best ArgCSR setting for only \stat{p2-d1-best-count-fp64} of \stat{p2-matrices} matrices in fp64. Fixing it
costs \stat{p2-d1-slowdown}\x{} on average against the best value per matrix, and up to \stat{p2-d1-worst}\x{}
(\texttt{\stat{p2-d1-worst-matrix}}). This is the parameter choice RASP automates.

\textbf{Paper~2's headline does not hold against 2026 cuSPARSE.} At its default, ArgCSR reaches
\stat{p2-argcsr-default-fp32}\x{} (fp32) and \stat{p2-argcsr-default-fp64}\x{} (fp64) of cuSPARSE \texttt{ALG1}
(geometric mean), and \stat{p2-argcsr-tuned-fp32}\x{} and \stat{p2-argcsr-tuned-fp64}\x{} when tuned per matrix. The
exceptions sit in the family the paper highlighted: on \texttt{ncvxqp9} (GHS\_indef) ArgCSR is
\stat{p2-ncvxqp9-argcsr-vs-cusparse}\x{} as fast as cuSPARSE, and tuned ArgCSR is the overall winner on \texttt{torso2}
(\cref{tab:named}). The vendor library has moved from the scalar CSR kernel of 2012 to load-balanced kernels, so the
qualitative trends of paper~2 survive while its absolute advantage does not.

\subsection{No format wins everywhere}
\begin{figure}[t]
\centering\rfig[0.95\linewidth]{p2_perf_profile}
\caption{Performance profiles over the \stat{p2-matrices} training matrices: the fraction of matrices on which each
format is within a factor $\tau$ of the fastest of all configurations.}\label{fig:profile}
\end{figure}

The fastest of all configurations (the \emph{oracle}) is \stat{p2-oracle-fp32}\x{} faster than cuSPARSE \texttt{ALG1}
in fp32 and \stat{p2-oracle-fp64}\x{} in fp64 (geometric mean; up to \stat{p2-oracle-fp32-max}\x{} and
\stat{p2-oracle-fp64-max}\x). cuSPARSE itself is the oracle on \stat{p2-cusparse-wins-fp32} (fp32) and
\stat{p2-cusparse-wins-fp64} (fp64) matrices; CSR-vector with a small $v$ wins most of the others (\cref{tab:wins},
\cref{fig:profile}). The headroom shrinks with size: \stat{p2-oracle-size-lt100k}\x{} below $10^5$ nonzeros,
\stat{p2-oracle-size-100k-1M}\x{} up to $10^6$ and \stat{p2-oracle-size-gt1M}\x{} above. Thread-per-row formats
collapse on power-law rows: on \texttt{ASIC\_100k}, a few rows with tens of thousands of nonzeros serialise on single
threads. This is why the selector needs the maximum row length.

Our two extensions behave as designed. One \emph{untuned} ArgCSR-local configuration is
\stat{p2-local-vs-default}\x{} faster than paper~2's default and reaches \stat{p2-local-vs-tuned} of ArgCSR tuned per
matrix over 18 configurations: choosing the chunk size locally recovers most of the tuning benefit. CSR-binned matches
per-matrix-tuned CSR-vector without tuning. Neither moves the oracle by more than about a percent, because modern
cuSPARSE is already near-optimal on large matrices. We therefore re-scoped \G2 from ``beat cuSPARSE'' to ``within
10\% of the oracle, and never lose on large matrices''.

\begin{table}[t]
\centering\small
\caption{Which family is fastest, per matrix and precision.}\label{tab:wins}
\stattable{p2-wins}
\end{table}

\subsection{G2--G4: the runtime selector on held-out matrices}\label{sec:eval-g2}
\begin{table}[t]
\centering\small
\caption{Steady-state selection on \stat{p3-heldout-samples} held-out samples (\stat{p3-heldout-matrices} matrices
$\times$ 2 precisions).}\label{tab:g2}
\stattable{p3-g2}
\end{table}

\begin{figure}[t]
\centering
\begin{minipage}[t]{0.49\linewidth}\centering\rfig{p3_oracle_gap}
\caption{Oracle gap on held-out samples: the CDF of $T_\text{selected}/T_\text{oracle}$.}\label{fig:gap}\end{minipage}\hfill
\begin{minipage}[t]{0.49\linewidth}\centering\rfig{p3_calls_sweep}
\caption{Plan cost plus $N$ calls, against always using cuSPARSE: break-even keeps small $N$ safe.}\label{fig:sweep}\end{minipage}
\end{figure}

\textbf{G2 is met on unseen matrices.} The decision tree runs within \stat{p3-g2-tree}\x{} of the held-out oracle
(target 1.10), within 10\% of it on \stat{p3-g2-tree-within10}\% of samples, and \stat{p3-g2-tree-vs-cusparse}\x{}
faster than always using cuSPARSE, which sits at \stat{p3-g2-cusparse}\x{} of the oracle (\cref{tab:g2},
\cref{fig:gap}). It is never more than 5\% slower than cuSPARSE on a matrix with at least $10^6$ nonzeros
(\stat{p3-g2-tree-slower-large} cases). Only \stat{p3-heldout-large} held-out samples are that large, so this half of
the goal rests mainly on cross-validation, where the gated tree had \stat{p3-cv-gated-large-slower} such losses.
Cross-validation predicted \stat{p3-cv-gated}; held-out performance is close to it. The hand-written rules (v1) are no
better than cuSPARSE (\stat{p3-g2-rules}\x{} of the oracle): thresholds fitted by hand do not transfer. The oracle
within the five-configuration portfolio (\stat{p3-g2-portfolio}\x) shows that the remaining headroom is in the
selector, not in the kernels.

\textbf{G3 is borderline.} Feature extraction takes a median of \stat{p3-g3-feature-us}\,\textmu s, which is
\stat{p3-g3-median} cuSPARSE SpMVs (90th percentile \stat{p3-g3-p90}). It is mostly a fixed launch, copy and
synchronisation cost under WSL: \stat{p3-g3-size-gt1M} SpMVs above $10^6$ nonzeros, but \stat{p3-g3-size-lt100k}
below $10^5$. That is why the feature gate skips extraction when few calls are expected.

\textbf{G4: break-even makes the selector safe at small $N$.} With break-even, the tree's plan plus $N$ calls runs at
\stat{p3-g4-tree-1}\x{} of always-cuSPARSE at $N=1$, and at $N=1$ \stat{p3-g4-n1-same} of \stat{p3-g4-n1-total} plans
do exactly the baseline's work (no features, same kernel); the remaining spread is the timing noise of a sub-millisecond
plan construction. Without break-even the same tree costs \stat{p3-g4-tree-nobe-1}\x{} at $N=1$ (worst
\stat{p3-g4-tree-nobe-1-worst}\x). At $N=1000$ the tree gains \stat{p3-g4-tree-1000}\x{} (\cref{fig:sweep}). Features
are extracted for \stat{p3-features-at-10}\% of plans at $N=10$, \stat{p3-features-at-100}\% at $N=100$ and
\stat{p3-features-at-1000}\% at $N=1000$.

\begin{table}[t]
\centering\small
\caption{Total time (plan + $N$ calls) vs always-cuSPARSE: geometric mean [worst case].}\label{tab:g4}
\stattable{p3-g4}
\end{table}

\subsection{G5: where should the operands live? (paper 1, Table 3)}\label{sec:eval-g5}
\begin{table}[t]
\centering\small
\caption{cuSPARSE SpGEMM with each operand in device (D), pinned zero-copy (P) or managed (M) memory: slowdown
against all-device. Our version of paper~1's Table~3.}\label{tab:placement}
\stattable{p4-placement}
\end{table}

\Cref{tab:placement} repeats paper~1's study at the level of the whole SpGEMM. With $B$ in pinned memory, $R\cdot A$
slows down \stat{p4-place-RxA-b}\x{} and $A\cdot P$ \stat{p4-place-AxP-b}\x{} (paper~1 on NVLink: 7--29\x). $A$ alone
in pinned memory costs \stat{p4-place-RxA-a}\x{} and \stat{p4-place-AxP-a}\x, and $C$ alone \stat{p4-place-RxA-c}\x{}
and \stat{p4-place-AxP-c}\x. Managed memory behaves exactly like pinned memory under WSL2 (all-managed vs all-pinned:
\stat{p4-managed-vs-pinned}\x). The flop-count bound overestimates $\mathrm{nnz}(C)$: the real compression is
\stat{p4-gamma-AxP} for $A\cdot P$ and \stat{p4-gamma-RxA} for $R\cdot A$, which is why RASP measures it instead of
sizing $C$ by the bound.
\todo{interpret the table once measured: which placements are tolerable on PCIe, and how this compares with the NVLink ratios.}

\subsection{G6: overlapped, adaptive out-of-core SpGEMM}\label{sec:eval-g6}
\begin{figure}[t]
\centering\rfig[\linewidth]{p4_scaling}
\caption{Out-of-core SpGEMM throughput (pinned inputs in, pinned output out). Left: fixed 512\,MiB budget, grid size
swept; right: fixed grid, budget swept. In the style of paper~1, Figs.~12--13.}\label{fig:scaling}
\end{figure}

\begin{table}[t]
\centering\small
\caption{Ablation: speedup over Algorithm~4 (no overlap) on all configurations where the problem exceeds the budget
(geometric mean).}\label{tab:ablation}
\stattable{p4-ablation}
\end{table}

When the problem exceeds the budget, RASP is \stat{p4-g6-AxP}\x{} faster than Algorithm~4 on $A\cdot P$ and
\stat{p4-g6-RxA}\x{} on $R\cdot A$ (target 1.3; verdicts: \stat{p4-g6-AxP-verdict}, \stat{p4-g6-RxA-verdict}).
\Cref{tab:ablation} attributes the gain step by step: double buffering alone with Algorithm~4's partitions
(\stat{p4-speedup-alg4-overlap-AxP}\x{} on $A\cdot P$), then the cost model's memory split and block sizes
(\stat{p4-speedup-rasp-static-AxP}\x), then calibration and re-planning (\stat{p4-speedup-rasp-AxP}\x). Zero-copy reaches
\stat{p4-speedup-zerocopy-AxP}\x{} of Algorithm~4. The measured parameters were a kernel rate of
\stat{p4-rate-gflops}\,GFLOP/s, a compression of \stat{p4-measured-gamma} and \stat{p4-measured-ws-per-flop} bytes of
cuSPARSE workspace per multiply-add. The cost model's predicted execution time was within a factor
\stat{p4-model-accuracy} of the measured one (median), and estimation plus planning took \stat{p4-overhead-pct}\% of
the total. Out of core, RASP reaches \stat{p4-frac-of-device} of the device-only kernel's throughput: what PCIe leaves
on the table. \Cref{fig:timeline} shows both schedules for the most out-of-core configuration.
\todo{interpret after the runs: where each step of the ablation helps, how often re-planning fired, behaviour as the
problem/budget ratio grows, and the in-core comparison when everything fits.}

\begin{figure}[t]
\centering\rfig[\linewidth]{p4_timeline}
\caption{Copy engine and compute stream over time for Algorithm~4 (sequential) and RASP (pipelined), from per-tile CUDA
events. \texttt{figures/nsys\_overlap.png} shows the same run in Nsight Systems.}\label{fig:timeline}
\end{figure}

\subsection{G7: end-to-end applications}\label{sec:eval-g7}
\paragraph{Conjugate gradients.} A solver knows its iteration cap, and that is the natural \texttt{expected\_calls}
hint for break-even. Over \stat{p5-cg-matrices} SPD matrices, Jacobi-preconditioned CG with \texttt{SpmvPlan} given that
hint takes \stat{p5-cg-tree-geomean}\x{} the time of the same solver with cuSPARSE (worst \stat{p5-cg-tree-worst}\x,
best \stat{p5-cg-tree-best}\x), and \stat{p5-cg-tree-exact-geomean}\x{} when it knows the exact iteration count
(\cref{tab:cg}, \cref{fig:cg}). SpMV is only \stat{p5-cg-spmv-share}\% of a solve (median); vector updates, reductions
and two host synchronisations per iteration take the rest, which caps what any SpMV selector can gain end to end.
\stat{p5-cg-failures} solves failed the CPU residual check.
\todo{interpret: when the iteration cap overstates the work, does break-even over-invest? Compare tree vs tree\_exact.}

\begin{table}[t]
\centering\small
\caption{CG time-to-solution (plan + solve) against the same solver with cuSPARSE.}\label{tab:cg}
\stattable{p5-cg}
\end{table}

\begin{figure}[t]
\centering\rfig[0.95\linewidth]{p5_cg}
\caption{CG speedup per matrix (sorted by nnz) for the three ways of giving SpmvPlan the call count.}\label{fig:cg}
\end{figure}

\paragraph{AMG set-up.} The Galerkin products $A_{l+1} = R_l(A_lP_l)$ of a Laplace3D hierarchy run through
\texttt{SpgemmContext} under a fast-memory budget (\cref{tab:amg}). Every method produces the same coarse operators
and satisfies the Galerkin row-sum identity on every level.
\todo{interpret after the runs: which levels run out of core, and where the set-up time goes.}

\begin{table}[t]
\centering\small
\caption{AMG set-up for Laplace3D $n=192$ under a 256\,MiB budget (median of 3).}\label{tab:amg}
\stattable{p5-amg-192}
\end{table}

\subsection{Kernel profiles}
\begin{table}[t]
\centering\small
\caption{Nsight Compute metrics of representative kernels (ncu locks clocks, so these are not timings).}\label{tab:ncu}
\resizebox{\linewidth}{!}{\stattable{p5-ncu}}
\end{table}
\Cref{tab:ncu} lists hardware counters for kernels behind the results above. CSR-vector on \texttt{af\_shell1}
reaches \stat{p5-ncu-spmv-csr-vector-af-shell1-dram}\% of peak DRAM throughput; ArgCSR on \texttt{rajat23} reaches
\stat{p5-ncu-spmv-argcsr-D32-rajat23-dram}\% with chunk size 32 and \stat{p5-ncu-spmv-argcsr-D1-rajat23-dram}\% with
chunk size~1 (\cref{sec:eval-g1} gives the padding of each).
\todo{interpret once profiled: is CSR-vector bandwidth-bound, does ArgCSR D=32 spend its traffic on padding, and what
limits the partial-product summation kernels (stall reasons)?}

\section{A multi-core CPU baseline: the OpenMP version}\label{sec:omp}

Comparing one CPU core with a whole GPU is an unfair contest, and a reader can fairly ask whether the GPU wins
only because it is parallel. We therefore parallelised the sequential program (\texttt{seq/rasp\_seq.cpp}) with OpenMP, changing no algorithm, so
that all CPU cores can be compared with the GPU. The sequential, OpenMP and CUDA versions read the same inputs and
print the same phase breakdown. All numbers here come from one laptop CPU (an Intel Core i5-1235U: two
performance cores with two hardware threads each and eight efficiency cores, 12 hardware threads, under WSL2
Ubuntu, g++ 15 with \texttt{-O2 -march=native}), with the sequential baseline measured on the same machine.
Threads are pinned to separate cores (\texttt{OMP\_PROC\_BIND=close}, \texttt{OMP\_PLACES=cores}).

\subsection{What was parallelised}
\begin{table}[t]
\centering\footnotesize
\caption{Kernels of \texttt{seq/rasp\_omp.cpp} and how each is parallelised.}\label{tab:omp-kernels}
\begin{tabular}{@{}p{2.1cm}p{5.8cm}@{}}
\toprule
kernel & OpenMP \\
\midrule
SpMV CSR, SELL & \texttt{parallel for} over rows (CSR) or slices (SELL); loop schedule chosen at run time (\texttt{--schedule}) \\
SpMV ELL & one contiguous block of rows per thread, each running the original column-by-column loops (no fork/join per column) \\
\texttt{dot} & \texttt{reduction(+)}; the summation order changes, so CG may differ from the sequential run in rounding only \\
\texttt{axpy}, Jacobi, $p$-update, CG diagonal & \texttt{parallel for schedule(static)} \\
SpGEMM flop estimate & \texttt{parallel for} with a \texttt{reduction(+)} \\
SpGEMM symbolic pass & per-thread marker array; each row stores only its count, the prefix sum is sequential ($O(\text{rows})$) \\
SpGEMM numeric pass & per-thread accumulator, marker and list of touched columns; row $i$ writes only its own slice $[rp_i, rp_{i+1})$ of $C$, so there are no conflicts \\
\bottomrule
\end{tabular}
\end{table}

Rows have very different lengths, which makes the loop schedule matter for SpMV on irregular matrices: this is the
CPU version of the load-balancing problem ArgCSR addresses on the GPU. The SpGEMM loops use
\texttt{schedule(dynamic, 256)} for the same reason. File parsing and format conversion stay sequential, because
the comparisons do not time them.

\subsection{Correctness}
\texttt{rasp\_omp --verify} keeps the original sequential kernels and compares against them on the same inputs, at
1, 2 and all threads. SpMV in every format must agree to a relative error of $10^{-12}$; CG must converge to the
same tolerance within two iterations of the sequential count; and every AMG coarse matrix must have exactly the
sequential \texttt{row\_ptr} and \texttt{col}, with values equal to $10^{-12}$. All checks pass (in practice SpMV
CSR matches exactly, the other formats differ by about $10^{-16}$, and the CG and AMG results are identical).

\subsection{Scaling}
\begin{figure}[t]
\centering\rfig[0.62\linewidth]{omp_scaling}
\caption{Speedup over the sequential run against thread count (fastest of two sweeps). The dashed line is the
ratio of all-thread to single-thread streaming bandwidth (\stat{omp-bw-1}\,GB/s to \stat{omp-bw-all}\,GB/s,
\stat{omp-bw-ceiling}\x), the expected ceiling for memory-bound kernels.}\label{fig:omp-scaling}
\end{figure}

\begin{table}[t]
\centering\small
\caption{Speedup over the sequential version on the same machine (fastest of two sweeps). SpMV is the geometric
mean over six matrices; CG solves a $128^3$ Laplacian; AMG is the SpGEMM time of a $160^3$ Laplace hierarchy.}\label{tab:omp-scaling}
\stattable{omp-scaling}
\end{table}

\Cref{tab:omp-scaling} and \cref{fig:omp-scaling} show the result. Speedups stop at about \stat{omp-spmv-best-speedup}\x{} (SpMV, \stat{omp-spmv-best-threads}
threads), \stat{omp-cg-best-speedup}\x{} (CG, \stat{omp-cg-best-threads} threads) and \stat{omp-amg-best-speedup}\x{} (AMG SpGEMM,
\stat{omp-amg-best-threads} threads), against a hardware limit of 12. Within one sweep, where no baseline drift enters, the
best gains over the same code on one thread are \stat{omp-spmv-best-speedup-vs1}\x, \stat{omp-cg-best-speedup-vs1}\x{} and
\stat{omp-amg-best-speedup-vs1}\x. The OpenMP version on one thread runs as fast as the sequential version (back-to-back runs agree
within measurement noise), so the OpenMP overhead itself is negligible.

\textbf{Memory bandwidth is the ceiling for SpMV and CG.} These kernels do two flops per twelve bytes read, so they
are limited by memory bandwidth, and a CPU saturates it with few threads: one thread streams \stat{omp-bw-1}\,GB/s,
all \stat{omp-bw-threads} threads \stat{omp-bw-all}\,GB/s. CG, whose time is
dominated by SpMV and by streaming vector updates, flattens at about this ratio. SpMV reaches slightly above it on some matrices, plausibly because part of their
working set stays in cache. For comparison the RTX~4050 reads \stat{pf-device-read}\,GB/s: the GPU's advantage on these kernels
is mainly bandwidth, which no number of CPU cores supplies.

\textbf{The loop schedule is not resolved at this noise level.} \Cref{tab:omp-sched} compares the loop schedules at the highest thread
count. We would expect \texttt{dynamic} to help on matrices with a few very long rows (\texttt{rajat23}), as the
GPU's row grouping does; at this thread count the three schedules are dominated by run-to-run noise, so we draw no
conclusion about their ranking.

\begin{table}[t]
\centering\small
\caption{SpMV CSR with each loop schedule at the highest thread count (geometric mean over six matrices).}\label{tab:omp-sched}
\stattable{omp-schedule}
\end{table}

\textbf{SpGEMM is held back by sequential system time, not by memory bandwidth.} We expected the AMG set-up to scale
further than SpMV, since it is hash and accumulator work. It does not. In the numeric pass, \texttt{std::vector::resize}
zero-fills the output arrays on a single thread, and the kernel page faults that this causes are charged as system time:
\stat{omp-amg-numeric-sys-s-t1}\,s (\stat{omp-amg-numeric-sys-pct-t1}\% of the pass) on one thread
and \stat{omp-amg-numeric-sys-s-best}\,s (\stat{omp-amg-numeric-sys-pct-best}\% of the pass) at the best thread count. That time does not shrink with the thread count, so
Amdahl's law caps the speedup. The prefix sum of the symbolic pass is sequential too, but it is negligible next to this.

\textbf{Limits of this measurement.} The i5-1235U has two kinds of core and the sequential run may land on a
performance core with a high turbo clock, which makes parallel efficiency look lower than it is. The 12-thread
results use hyperthread siblings and efficiency cores and varied much more between the two sweeps than the
1--8 thread results, so we do not rely on them; sequential-baseline drift between runs is about $\pm$10\%.
WSL2 also inflates system time. The same sweep runs on the Kaggle and RTX~4050 machines, so every
comparison is against a sequential and a GPU run on the same hardware.

\subsection{OpenMP against the GPU}
Even with every core busy, the CPU version is limited to \stat{omp-cg-best-speedup}--\stat{omp-spmv-best-speedup}\x{} over one core on its memory-bound kernels,
because its bandwidth rises only \stat{omp-bw-ceiling}\x{} with all threads, while the GPU has several times the bandwidth.
\Cref{tab:same-machine} compares the GPU with the CPU on the same machines.


\end{document}
